\section{Related Work}
\label{sec:related_work}

\subsection{Weight Space Learning}

Neural network weights exist in high-dimensional spaces with rich geometric structure. Linear mode connectivity demonstrates that trained networks can be connected via low-loss paths in weight space~\cite{frankle2020linear}, suggesting meaningful structure beyond random initialization. This motivates learning representations that capture weight space semantics.

\textbf{Hyper-Representations}~\cite{schurholt2022hyper} train variational autoencoders on collections of neural network weights, demonstrating that latent spaces capture training dynamics and model relationships. However, evaluation focused on weight generation rather than property prediction, and comparisons used different datasets and metrics across papers.

\textbf{Statistical approaches} (StatNN and variants) compute simple statistics (mean, variance) over weights for property prediction. While computationally efficient, these methods ignore network structure entirely, treating all parameters as an unordered bag.

\subsection{Permutation Symmetries in Weight Space}

Neural networks exhibit permutation symmetries: reordering hidden neurons with corresponding weight adjustments preserves function. This symmetry complicates weight space learning---functionally identical networks may have distant representations.

\textbf{Git Re-Basin}~\cite{ainsworth2022git} addresses this through weight alignment, finding permutations that minimize distance between networks. Originally developed for model merging, alignment could preprocess weights before embedding. The computational cost scales quadratically with model count, creating challenges at scale.

\textbf{Neural Functional Transformers (NFN)}~\cite{zhou2024neural} take an architectural approach, designing layers that are permutation-equivariant by construction. This avoids alignment preprocessing but requires specialized architectures. NFN achieves strong results on weight processing tasks but has not been systematically compared to simpler alternatives for property prediction.

\subsection{Model Zoos and Benchmarks}

Large collections of pretrained models with metadata enable systematic evaluation. \textbf{Model Zoos}~\cite{schurholt2022model} provide thousands of checkpoints with ground-truth accuracy labels, establishing infrastructure for weight embedding benchmarks. However, existing evaluations compare methods on different zoo subsets with incompatible protocols.

\subsection{Gap in Existing Work}

Prior work evaluates embedding methods in isolation on different datasets with different metrics. We cannot determine which structural biases matter most for property prediction. This paper provides the missing systematic comparison through a controlled ablation study.
